\subsection{Compactness criterion for tight convergence}

DEFINITION 2. --- Let T be a topological space, and let H be a subset of \(\mathcal{M}^{\mathrm{b}}(\mathrm{T})\) ; one says that H satisfies Prokhorov's condition if

\begin{enumerate}
\def\labelenumi{\alph{enumi})}
\item
  \(\sup_{\mu \in \mathrm{H}}|\mu |(1) < +\infty ;\)
\item
  for every number \(\varepsilon > 0\) , there exists a compact subset \(\mathbf{K}_{\varepsilon}\) of \(\mathbf{T}\) such that
\end{enumerate}

\[
| \mu | (\mathrm{T} - \mathrm{K} _ {\varepsilon}) \leqslant \varepsilon \quad \text { for every measure } \mu \in \mathrm{H}.\tag{6}
\]

It can be shown that if T is completely regular, the set of conditions a) and b) is equivalent to the following condition: there exists a real function \(f \geqslant 1\) on T, such that the set of points t of T satisfying \(f(t) \leqslant c\) is compact for every \(c \in R_{+}\) (which in particular implies that f is lower semi-continuous), and such that \(\sup_{\mu \in \mathbf{H}} |\mu|(f) < +\infty\) . Moreover, when \(\mathbf{T}\) is locally compact, one obtains an equivalent statement by requiring \(f\) to be continuous (cf.~Exer. 10).

PROPOSITION 11. --- Let T be a completely regular space, and let H be a subset of \(\mathcal{M}^{b}(T)\) that satisfies Prokhorov's condition; then its closure \(\overline{H}\) in \(\mathcal{M}^{b}(T)\) satisfies Prokhorov's condition.

For, the functions \(\mu \mapsto |\mu|^{\bullet}(1)\) , \(\mu \mapsto |\mu|^{\bullet}(\mathrm{T} - \mathrm{K}_{\varepsilon})\) are lower semicontinuous on \(\mathcal{M}^b (\mathrm{T})\) by Prop. 6 of No.~3.

The interest of Prokhorov's condition comes from the following theorem, whose converse will be studied later on (Th. 2).

THEOREM 1 (Prokhorov). --- Let T be a completely regular space, and let H be a subset of \(\mathcal{M}^{b}(\mathrm{T})\) that satisfies Prokhorov's condition; then H is relatively compact in \(\mathcal{M}^{b}(\mathrm{T})\) for the tight topology.

We can suppose that T is a subspace of a compact space X; let i be the canonical injection of T into X. We can on the other hand suppose that H is closed in \(\mathcal{M}^{b}(T)\) , by Prop. 11. It will then suffice to show that every ultrafilter U on H converges in \(\mathcal{M}^{b}(T)\) .

We shall begin with the case that \(\mathrm{H}\subset\mathcal{M}_{+}^{b}(\mathrm{T})\) . The total masses of the measures \(\mu\in H\) being bounded by hypothesis, \(i(\mu)\) converges vaguely with respect to U, in \(\mathcal{M}_{+}(\mathrm{X})\) , to a measure \(\nu\in\mathcal{M}_{+}(\mathrm{X})\) (Ch. III, §1, No.~9, Cor. 2 of Prop. 15); by Prop. 8 of No.~3, everything comes down to proving that \(\nu\) is concentrated on T. Now, let \(\varepsilon\) be a number \textgreater0, and let \(K_{\varepsilon}\) be a compact subset of T satisfying the formula (6). Since \(X-K_{\varepsilon}\) is open in X, we have, by Prop. 6 of No.~3 applied in X, the inequalities

\[
\begin{array}{r l} \nu^ {\bullet} (\mathrm{X} - \mathrm{T}) & \leqslant \nu^ {\bullet} (\mathrm{X} - \mathrm{K} _ {\varepsilon}) \leqslant \liminf _ {\mu , \mathfrak {U}} i (\mu) ^ {\bullet} (\mathrm{X} - \mathrm{K} _ {\varepsilon}) \\ & = \liminf _ {\mu , \mathfrak {U}} \mu^ {\bullet} (\mathrm{T} - \mathrm{K} _ {\varepsilon}) \leqslant \varepsilon ; \end{array}
\]

since \(\varepsilon > 0\) is arbitrary, the theorem is established in this special case.

Let us pass to the general case; for every measure \(\mu\) on T, set

\[
a _ {1} (\mu) = \mathcal {R} (\mu) ^ {+}, a _ {2} (\mu) = \mathcal {R} (\mu) ^ {-}, a _ {3} (\mu) = \mathcal {I} (\mu) ^ {+}, a _ {4} (\mu) = \mathcal {I} (\mu) ^ {-};
\]

since \(\mu = a_{1}(\mu) - a_{2}(\mu) + ia_{3}(\mu) - ia_{4}(\mu)\) , it will suffice to show that the mappings \(a_{j}\) ( \(j = 1,2,3,4\) ) converge tightly with respect to \(\mathfrak{U}\) . But the set \(\mathbf{H}_j\) of measures \(a_{j}(\mu)\) , where \(\mu\) runs over \(\mathbf{H}\) , satisfies Prokhorov's condition by virtue of the relation \(|a_j(\mu)| \leqslant |\mu|\) , and is contained in \(\mathcal{M}_+^b(\mathrm{T})\) ; it is therefore relatively compact in \(\mathcal{M}_+^b(\mathrm{T})\) by the special case, and the theorem then follows at once.

COROLLARY. --- Let T be a subspace of a completely regular space X, and let H be a subset of \(\mathcal{M}^{b}(T)\) that satisfies Prokhorov's condition. If i denotes the canonical injection of T into X, then the restriction to H of the mapping \(\mu\mapsto i(\mu)\) of \(\mathcal{M}^{b}(T)\) into \(\mathcal{M}^{b}(X)\) is a homeomorphism of H onto its image.

It suffices to treat the case that H is closed (Prop. 11), hence compact; the conclusion then follows from the fact that \(\mu \mapsto i(\mu)\) is continuous and injective.

Recall that this result is also valid for an arbitrary subset of \(\mathcal{M}_{+}^{b}(\mathrm{T})\) (No.~3, Prop. 8).

THEOREM 2. --- Let T be a locally compact space, or a Polish space, and let H be a relatively compact subset of \(\mathcal{M}_{+}^{b}(T)\) ; then H satisfies Prokhorov's condition.

We may restrict ourselves to the case that H is closed, hence compact. The total masses of the measures \(\mu \in H\) are obviously bounded, because the mapping \(\mu \mapsto \mu(1)\) is continuous, and everything comes down to proving the assertion b) of Def. 2.

Suppose first that T is locally compact. Let \(\varepsilon\) be a number \textgreater0. Let us associate to every measure \(\mu\in H\) a compact set \(K_{\mu}\) in T such that \(\mu^{\bullet}(T-K_{\varepsilon})<\varepsilon\) , then a relatively compact open neighborhood \(U_{\mu}\) of \(K_{\mu}\) . The function \(\lambda\mapsto\lambda^{\bullet}(T-U_{\mu})\) being upper semi-continuous on \(\mathcal{M}_{+}^{b}(T)\) (No.~3, Prop. 6), the set \(V^{\mu}\) of measures \(\lambda\in H\) such that \(\lambda^{\bullet}(T-U_{\mu})<\varepsilon\) is a neighborhood of \(\mu\) in H. Therefore there exists a finite subset \(H'\) of H such that the sets \(V^{\mu}\) ( \(\mu\in H'\) ) cover H. Denoting by K the compact set \(\bigcup_{\mu\in H'}\overline{U}_{\mu}\) , we have \(\lambda^{\bullet}(T-K)<\varepsilon\) for all \(\lambda\in H\) .

Suppose next that T is Polish. We do not restrict the generality by assuming that T is the intersection of a decreasing sequence \((\mathrm{T}_{p})_{p\geqslant1}\) of open subsets of a compact space X (GT, IX, §6, No.~1, Cor. 1 of Th. 1). Let \(i_{p}\) be the injection of T into \(T_{p}\) , and let \(H_{p}\) be the set of measures of the form \(i_{p}(\lambda)\) for \(\lambda\in H\) ; since \(H_{p}\) is compact in \(\mathcal{M}_{+}^{b}(T_{p})\) , it follows that there exists a compact set \(K_{p}\subset T_{p}\) such that \(\nu^{\bullet}(T_{p}-K_{p})\leqslant\varepsilon2^{-p}\) for every measure \(\nu\in H_{p}\) , by the preceding result applied to the locally compact space \(T_{p}\) . Therefore also \(\nu^{\bullet}\big(\mathrm{T}-(\mathrm{T}\cap\mathrm{K}_{p})\big)\leqslant\varepsilon2^{-p}\) , and finally \(\lambda^{\bullet}\big(\mathrm{T}-(\mathrm{T}\cap\mathrm{K}_{p})\big)\leqslant\varepsilon2^{-p}\) for every measure \(\lambda\in H\) . Now set \(K=\bigcap_{p}K_{p}\) ; the set K is compact and is contained in T, and, for every measure \(\lambda\in H\) , we have \(\lambda^{\bullet}(T-K)\leqslant\sum_{p}\lambda^{\bullet}\big(\mathrm{T}-(\mathrm{T}\cap\mathrm{K}_{p})\big)\leqslant\sum_{p}\varepsilon2^{-p}=\varepsilon\) . Prokhorov's condition is thus verified.

